\documentclass[11pt]{beamer}
\usepackage[utf8]{inputenc}
\usepackage[T1]{fontenc}
\usepackage{lmodern}
\usepackage[french]{babel}
\usepackage{quiver}
\usetheme{Madrid}
\begin{document}
	\author[Maxime LOUVET]{Maxime LOUVET}
	\title[Audition EDSTS]{Cryptanalyse des systèmes à base de courbes elliptiques via isogénies}
	%\subtitle{test}
	%\logo{}
	%\institute{}
	\date[17/06/2025]{17/06/2025}
	%\subject{}
	%\setbeamercovered{transparent}
	%\setbeamertemplate{navigation symbols}{}
	\begin{frame}[plain]
		\maketitle
	\end{frame}
	
	\begin{frame}
		\frametitle{Profil académique}
		\begin{block}{Mon parcours universitaire}
			\begin{itemize}
				\item<1-> 2017-2020 \textbf{CPGE filière MPSI/MP*},
				Lycée Faidherbe, Lille.
				
				\item<2-> 2020-2023 \textbf{L3, Master Formation à l'Enseignement Supérieur de Mathématiques},
				Université Paris-Saclay, Orsay.
				
				\item<3-> 2024/2025 \textbf{M2 Algèbre Appliquée},
				Université Versailles Saint-Quentin.
				
			\end{itemize}
		\end{block}
	\end{frame}
	
	\begin{frame}
		\frametitle{Expérience professionnelle} 
		\begin{block}{Mon parcours universitaire}
			\begin{itemize}
				
				\item<1-> 09/2023-02/2024 \textbf{Professeur agrégé de mathématiques}
				Lycée Jules Verne, Limours.
				
				\item<2-> 03/2024-Aujourd'hui \textbf{Professeur particulier chez Acadomia}, niveau lycée, licence et CPGE.
				
				\item<3-> 08/2024-02/2025 \textbf{Khôlleur}, lycée Lakanal, Sceaux.
				
				\item<4-> 02/2025-Aujourd'hui \textbf{Stage de M2 au laboratoire MIS}, Amiens.
				
			\end{itemize}
		\end{block}
	\end{frame}
	
	\begin{frame}
		\frametitle{Échange de clef en cryptographie} 
		Alice et Bob veulent créer un secret $K$ à partir d'une donnée $P$ publique. 
		\begin{center}
			Soit $(G,+)$ un groupe, et $P\in(G,+)$.
			
			\includegraphics[scale = 0.4]{DH}
		
		\begin{block}{Problème du logarithme discret (DLP)}
			Étant donné $P$ et $Q$ deux éléments d'un groupe $G$, trouver $k$ tel que $Q = k\cdot P$.
		\end{block}
	\end{center}
	\end{frame}
	
	\begin{frame}
		\frametitle{Groupe et géométrie} 
		\begin{center}
			\includegraphics[scale = 0.4]{cercle1}<1>
			\includegraphics[scale = 0.4]{cercle2}<2>
			\includegraphics[scale = 0.4]{cercle3}<3>
			\includegraphics[scale = 0.4]{cercle4}<4>
		\end{center}
		
	\end{frame}
	
	\begin{frame}
		\frametitle{Courbes elliptiques} 
		\begin{itemize}
			\item Problème DLP sur le cercle : Facile.
			\item Droites, ellipses, hyperboles... : courbe de "genre" $0$ donc "simple".
		\begin{block}{Courbes Elliptiques}
			Les courbes elliptiques sont les courbes de genre $1$.
		\end{block}
		
		\item Problème DLP sur courbes elliptiques : Difficile. 
		\end{itemize}
		\begin{block}{Projet réalisé}
			\begin{itemize}
				\item 2023 \textbf{Baby step Giant step} Problème DLP sur groupe quelconque.
				\item 2025 \textbf{Méthode GLV} Calcul de $k\cdot A$ sur courbe elliptique.
			\end{itemize}
		\end{block}
	\end{frame}
	
	\begin{frame}
		\frametitle{Stage actuel} 
		\framesubtitle{Cryptanalyse et isogénies} 
		
		Si $E$ et $E'$ sont deux courbes elliptiques, soit $\phi$ une application qui à un point $P\in E$ associe un point $\phi(P)\in E'$.
		
		\begin{block}{Isogénie entre courbes elliptiques}
			$\phi$ est une isogénie si elle vérifie
			\begin{equation*}
				\phi(P+Q) = \phi(P) + \phi(Q).
			\end{equation*}
		\end{block}
		
		\begin{block}{Lien avec le problème DLP }
			Si $\phi$ est une isogénie, alors $Q = k\cdot P\iff \phi(Q) = k\cdot\phi(P)$
		\end{block}
		
	\end{frame}
	
	\begin{frame}
		\frametitle{Stage actuel} 
		\framesubtitle{Calculer une isogénie}
		\begin{block}{Sujet de stage}
			Trouver $\phi$ isogénie entre $E$ et $E'$ ? Calculer $\phi(P)$ lorsque $P\in E$?
		\end{block}
		\begin{itemize}
			\item Difficulté mesurée selon $q$ (taille des groupes) et $l$ ("degrés" de l'isogénie).
			\item Formule de Vélu : méthode générale, rapide selon $q$ mais lent selon $l$ : Permet de calculer les "petites" isogénies.
		\end{itemize}
		
		
		\begin{block}{Méthodes de calcul d'isogénies}
			\begin{itemize}
				\item Idée 1 : \textbf{Action du groupe de classes}(2008): rapide selon $l$, mais moyennement rapide selon $q$. 
				\item Idée 2 : \textbf{Lemme de Kani}(2022): rapide en théorie selon $l$ et $q$. 
			\end{itemize}
		\end{block}
	\end{frame}
	
	\begin{frame}
		\frametitle{Stage actuel} 
		\framesubtitle{Action du groupe de classe}
		\begin{block}{Action du groupe de classe}
			\begin{enumerate}
				\item On représente une isogénie par un objet $L$ et une classe $\tilde{L}$.
				\item $\tilde{L}$ s'écrit dans un groupe $\tilde{L} = \tilde{L_1}\cdots\tilde{L_n}$. 
				\item On calcule les isogénies associées à $L_1, \cdots, L_n$ avec Vélu.
				\item On les enchaîne pour calculer $\phi$.
			\end{enumerate}
		\end{block}
		\begin{itemize}
			\item $L_1,\cdots,L_n$ : petites isogénies, donc la formule de Vélu est rapide.
			\item Trouver $\tilde{L} = \tilde{L_1}\cdots\tilde{L_n}$ lent selon $q$ : Optimisation possible. 
		\end{itemize}
		
		\begin{block}{Contribution personnelle}
			\begin{itemize}
				\item Proposition d'optimisation pour l'étape $2$, implémentée dans Sage. 
				\item Comparaison avec l'optimisation initiale (Jao Soukharev 2010).
			\end{itemize}
		
		\end{block}
	
	\end{frame}
	
	\begin{frame}
		\frametitle{Introduction à la thèse} 
		
		
	\end{frame}
	
\end{document}